%!TEX root = ../template.tex

%% ----------------------------------------------------------------------------
%%  NOVATHESIS — 2-MainMatter/chapter-build-py.tex
%% ----------------------------------------------------------------------------
%%
%% Version 7.10.7 (2026-03-25)
%% Copyright (C) 2004-26 by João M. Lourenço <joao.lourenco@fct.unl.pt>
%%
%% Tutorial chapter: the .Build/build.py build assistant.
%% Intended audience: developers and CI users who need to build NOVAthesis
%% documents programmatically for specific schools, document types, and
%% languages without manually editing configuration files.
%%
%% ----------------------------------------------------------------------------

\typeout{NT FILE 2-MainMatter/chapter-build-py.tex}%

\chapter{The \texttt{build.py} Build Assistant}
\label{app:build_py}

The file \verb!.Build/build.py! is a Python~3 script that automates building
\novathesis{} documents from the command line without touching any
configuration file by hand.  It is the primary tool used by the
CI pipeline and by maintainers generating sample PDFs for multiple
schools at once.


%%=======================================================================
\section{Overview}
\label{sec:buildpy_overview}
%%=======================================================================

The script performs four steps in sequence:

\begin{enumerate}
  \item \textbf{Resolve the school.}  Validate the \texttt{school\_id}
        argument and auto-detect the highest available degree if
        \texttt{--doctype} was not supplied.
  \item \textbf{Prepare a workspace.}  Depending on the options, either
        use the current directory in-place or create a temporary directory
        populated with symlinks to the project tree.
  \item \textbf{Patch configuration files.}  If \emph{patching mode} is
        active (see Section~\ref{sec:buildpy_modes}), the copies of
        \verb!0-Config/1_novathesis.tex! and \verb!0-Config/4_files.tex!
        in the workspace are modified to reflect the requested school,
        document type, language, and status.
  \item \textbf{Build and deliver.}  Run \texttt{make} inside the workspace,
        then copy the resulting \verb!template.pdf! to the output directory.
\end{enumerate}

\medskip\noindent
\textbf{When to use \texttt{build.py} vs.\ \texttt{make}.}  For everyday
writing, just run \texttt{make} (or your editor's build command).
Use \texttt{build.py} when you need to build for a \emph{different} school
or document type than the one currently set in \verb!0-Config/!, or when
building from a script or CI job.


%%=======================================================================
\section{Prerequisites}
\label{sec:buildpy_prereqs}
%%=======================================================================

\begin{itemize}
  \item Python~3.10 or later.
  \item A complete \TeX~Live (or MiK\TeX) installation with \texttt{latexmk}.
  \item The script must be run from the \textbf{project root} — the directory
        that contains \verb!template.tex! and \verb!0-Config/!.
\end{itemize}

\begin{bashcode}
# From the project root:
python3 .Build/build.py nova/fct
\end{bashcode}


%%=======================================================================
\section{Quick Start}
\label{sec:buildpy_quickstart}
%%=======================================================================

The only required argument is the \emph{school identifier}: a
\verb!university/faculty! path that mirrors the directory structure under
\verb!NOVAthesisFiles/Schools/!.

\begin{bashcode}
# Simplest call — build for NOVA FCT in user mode
python3 .Build/build.py nova/fct

# Build a PhD thesis in English, patching configs automatically
python3 .Build/build.py nova/fct --doctype phd --lang en --docstatus final

# Build cover pages only (fast preview)
python3 .Build/build.py nova/fct --cover

# Dry run — show what would happen without building
python3 .Build/build.py nova/fct --doctype msc --dry-run
\end{bashcode}

The separator between university and faculty may be \verb!/!, \verb!-!, or
\verb!.!; all are normalised to \verb!/! internally.  So \verb!nova-fct!
and \verb!nova/fct! are identical.


%%=======================================================================
\section{Build Modes}
\label{sec:buildpy_modes}
%%=======================================================================

The script operates in one of three modes, selected automatically from the
options you pass.

%%-----------------------------------------------------------------------
\subsection{User Mode}
\label{sec:buildpy_user_mode}
%%-----------------------------------------------------------------------

\emph{User mode} is active when \verb!--docstatus! is \texttt{keep} (the
default) and neither \verb!--cover! nor \verb!--force-school! is given.
In this mode:

\begin{itemize}
  \item The configuration files in \verb!0-Config/! are \emph{not modified}.
  \item The build runs in the \textbf{current directory} (no temporary
        workspace is created).
  \item The result is exactly what \texttt{make} would produce.
\end{itemize}

\begin{bashcode}
python3 .Build/build.py nova/fct         # user mode (default)
python3 .Build/build.py nova/fct --progress 2  # same, with real-time output
\end{bashcode}

%%-----------------------------------------------------------------------
\subsection{Patching Mode}
\label{sec:buildpy_patching_mode}
%%-----------------------------------------------------------------------

\emph{Patching mode} is activated whenever any of the following options are
present: \verb!--docstatus! (with a value other than \texttt{keep}),
\verb!--cover!, or \verb!--force-school!.

In patching mode the script:

\begin{enumerate}
  \item Automatically creates a \textbf{temporary workspace} (unless you
        specify a directory with \verb!--build-dir!).
  \item Populates it with symlinks pointing back to the project tree, then
        copies (not symlinks) the \verb!0-Config/! files so they can be
        safely modified.
  \item Patches the copies to set the requested school, doctype, language,
        and status.
  \item Builds inside the workspace; the original \verb!0-Config/! files are
        never touched.
\end{enumerate}

\begin{bashcode}
# Patch school, doctype, language and status — builds in a temp directory
python3 .Build/build.py nova/fct \
    --doctype phd \
    --lang en \
    --docstatus final \
    --rename-pdf \
    --output-dir build/
\end{bashcode}

%%-----------------------------------------------------------------------
\subsection{Cover-Only Mode}
\label{sec:buildpy_cover_mode}
%%-----------------------------------------------------------------------

The \verb!--cover! flag builds \emph{only the cover pages}, skipping all
chapters and bibliography.  This is much faster than a full build and is
useful for previewing cover layouts.

\begin{itemize}
  \item Implies \verb!--docstatus final! automatically.
  \item Comments out all \verb!\ntaddfile! entries in
        \verb!0-Config/4_files.tex! in the workspace.
  \item Disables the copyright page (\verb!\ntsetup{print/copyright=false}!).
  \item Expected line count is reduced to \(\approx 2400\) for a more
        accurate progress bar.
\end{itemize}

\begin{bashcode}
python3 .Build/build.py nova/fct --cover
python3 .Build/build.py nova/fct --cover --lang pt --progress 2
\end{bashcode}


%%=======================================================================
\section{Command-Line Reference}
\label{sec:buildpy_reference}
%%=======================================================================

%%-----------------------------------------------------------------------
\subsection{Positional Argument}
\label{sec:buildpy_positional}
%%-----------------------------------------------------------------------

\begin{description}
  \item[\texttt{school\_id}]
        The school identifier in \texttt{university/faculty} format.  For
        sub-department variants use a three-part path such as
        \texttt{nova/fct/cbbi}.  Hyphens and dots are accepted as separators
        and are normalised to slashes.

        \textbf{Examples:}
        \verb!nova/fct!, \verb!nova/fcsh!, \verb!ulisboa/ist!,
        \verb!uminho/eeng!, \verb!nova-fct! (same as \verb!nova/fct!).
\end{description}

%%-----------------------------------------------------------------------
\subsection{Document Options}
\label{sec:buildpy_doc_opts}
%%-----------------------------------------------------------------------

\begin{description}
  \item[\texttt{-t}, \texttt{--doctype} \{phd | msc | bsc\}]
        Document type.  If omitted, the script consults
        \verb!.Build/schools.conf! and selects the highest degree available
        for the given school (usually \texttt{phd}).

  \item[\texttt{-s}, \texttt{--docstatus} \{working | provisional | final | keep\}]
        Document status written into \verb!\ntsetup{docstatus=...}!.
        \texttt{keep} (the default) leaves whatever value is already in
        \verb!0-Config/1_novathesis.tex! unchanged.  Any other value
        activates patching mode.

  \item[\texttt{-l}, \texttt{--lang} \textit{code}]
        Two- or three-letter \textsc{bcp-47} language code (default:
        \texttt{en}).  Sets \verb!\ntsetup{lang=...}!.  Examples:
        \texttt{en}, \texttt{pt}, \texttt{de}, \texttt{fr}, \texttt{uk},
        \texttt{gr}, \texttt{cat}.

  \item[\texttt{-{}-sdgs} \textit{n,n,\ldots}]
        Comma-separated list of \textsc{sdg} numbers to display on the cover
        (default: all 17).  Sets \verb!\ntsetup{print/sdgs/list=...}!.

  \item[\texttt{-c}, \texttt{--cover}]
        Build cover pages only.  Activates cover-only mode (see
        Section~\ref{sec:buildpy_cover_mode}).

  \item[\texttt{-f}, \texttt{--force-school}]
        Force the school value to be applied even in user mode.  Useful when
        the school in \verb!0-Config/1_novathesis.tex! differs from the one
        passed on the command line.
\end{description}

%%-----------------------------------------------------------------------
\subsection{Build Directory Options}
\label{sec:buildpy_bdir_opts}
%%-----------------------------------------------------------------------

\begin{description}
  \item[\texttt{-bdir} / \texttt{--build-dir} \textit{[path | - ]}]
        Controls where the build workspace lives:

        \begin{center}
        \renewcommand{\arraystretch}{1.25}
        \begin{ShadedTabular}{lp{9cm}}
          \toprule
          \textbf{Value} & \textbf{Effect} \\
          % \midrule
          Omitted           & Use current directory (user mode default) \\
          Present, no value & Create a new \verb!tempfile! directory \\
          \texttt{-}        & Reuse the directory cached in
                              \verb!.Build/.keep-dir.txt! from a previous
                              run; create a new temp dir if none is cached \\
          \textit{path}     & Use (or create) the specified directory \\
          \bottomrule
        \end{ShadedTabular}
        \end{center}

  \item[\texttt{-k}, \texttt{--keep-bdir}]
        Do not delete the temporary workspace after a successful build.
        Implied when using \texttt{-bdir -} (cached directory mode).

  \item[\texttt{-o}, \texttt{--output-dir} \textit{path}]
        Directory where the finished PDF is copied (default: current
        directory).  Created automatically if it does not exist.

  \item[\texttt{-r}, \texttt{--rename-pdf}]
        Rename the output PDF from the generic \verb!template.pdf! to
        \verb!<university>-<faculty>-<doctype>-<lang>.pdf!.  For example:
        \verb!nova-fct-phd-en.pdf!.
\end{description}

%%-----------------------------------------------------------------------
\subsection{Build Engine Options}
\label{sec:buildpy_engine_opts}
%%-----------------------------------------------------------------------

\begin{description}
  \item[\texttt{-p}, \texttt{--processor} \textit{name}]
        The \texttt{make} target used to compile \hologo{LaTeX}
        (default: \texttt{lua}).  Pass \texttt{pdf} for \hologo{pdfLaTeX} or
        \texttt{xe} for \hologo{XeLaTeX}.

  \item[\texttt{-{}-progress} \{0 | 1 | 2\}]
        Output verbosity during compilation:
        \begin{itemize}
          \item \texttt{0} — silent (no output)
          \item \texttt{1} — animated progress bar (default)
          \item \texttt{2} — real-time \texttt{latexmk} output
        \end{itemize}

  \item[\texttt{-{}-lines} \textit{n}]
        Expected number of output lines for the progress-bar calculation.
        If omitted, the value from a previous run is read from
        \verb!.Build/.keep-lines.txt!; falls back to 4400 for a full build
        and 2400 for cover-only mode.
\end{description}

%%-----------------------------------------------------------------------
\subsection{Utility Options}
\label{sec:buildpy_util_opts}
%%-----------------------------------------------------------------------

\begin{description}
  \item[\texttt{-{}-dry-run}]
        Print what the script \emph{would} do — school, doctype, language,
        build directory, files to modify — without running the build.

  \item[\texttt{-v}, \texttt{--verbose}]
        Enable extra debug output from the script itself.
\end{description}


%%=======================================================================
\section{Doctype Auto-Detection}
\label{sec:buildpy_autodetect}
%%=======================================================================

When \verb!--doctype! is omitted, \texttt{build.py} reads
\verb!.Build/schools.conf! to determine the highest degree available for the
requested school.  Each line in that file lists a school identifier followed
by the degrees it supports, e.g.:

\begin{verbatim}
nova/fct    [phd, msc]  [en, pt]
uminho/ec   [phd, msc]  [en, pt]
ipl/isel    [msc]       [en, pt]
\end{verbatim}

The script picks the \emph{first} entry in the degree list (leftmost = highest).
If the school is not listed in \verb!schools.conf!, the fallback is
\texttt{phd}.  You can always override auto-detection with \verb!--doctype!.


%%=======================================================================
\section{Examples}
\label{sec:buildpy_examples}
%%=======================================================================

%%-----------------------------------------------------------------------
\subsection{CI Pipeline Build}
\label{sec:buildpy_ci}
%%-----------------------------------------------------------------------

The GitHub Actions workflow uses:

\begin{bashcode}
python3 .Build/build.py nova/fct \
    --build-dir=AUXDIR \
    --progress=2
\end{bashcode}

\verb!--progress 2! streams all \texttt{latexmk} output directly to the
CI log.  \verb!--build-dir AUXDIR! places all auxiliary files in a named
directory so they can be archived as a workflow artifact.

%%-----------------------------------------------------------------------
\subsection{Generating Multiple School PDFs}
\label{sec:buildpy_multi}
%%-----------------------------------------------------------------------

A shell loop can drive \texttt{build.py} to produce one PDF per school:

\begin{bashcode}
for school in nova/fct nova/fcsh ulisboa/ist uminho/eeng; do
    python3 .Build/build.py "$school"     \
        --doctype phd                     \
        --lang en                         \
        --docstatus final                 \
        --rename-pdf                      \
        --output-dir pdfs/
done
\end{bashcode}

%%-----------------------------------------------------------------------
\subsection{Iterative Cover Development}
\label{sec:buildpy_iterative}
%%-----------------------------------------------------------------------

While working on a new school cover, iterate quickly with:

\begin{bashcode}
# First pass — creates a temp directory and caches its path
python3 .Build/build.py nova/fct --cover --build-dir -

# Subsequent passes — reuse the same temp directory
python3 .Build/build.py nova/fct --cover --build-dir -
\end{bashcode}

Using \verb!--build-dir -! on the first call creates a temporary directory
and saves its path to \verb!.Build/.keep-dir.txt!.  All later calls with
\verb!--build-dir -! reuse that directory, skipping the symlink-creation
step and making repeated builds noticeably faster.

%%-----------------------------------------------------------------------
\subsection{Dry Run Before Committing}
\label{sec:buildpy_dryrun}
%%-----------------------------------------------------------------------

Before a large batch run, verify the parameters with \verb!--dry-run!:

\begin{bashcode}
python3 .Build/build.py nova/fct \
    --doctype msc \
    --lang pt \
    --docstatus provisional \
    --dry-run
\end{bashcode}

The script prints the resolved school, document type, language, build
directory, and the configuration files it would modify, then exits without
building.


%%=======================================================================
\section{Quick Reference}
\label{sec:buildpy_quickref}
%%=======================================================================

\begin{center}
\renewcommand{\arraystretch}{1.3}
\begin{ShadedTabular}{lp{8.5cm}}
  \toprule
  \textbf{Option} & \textbf{Effect} \\
  % \midrule
  \verb!school_id!               & Required. \texttt{university/faculty[/variant]} \\
  \verb!-t! / \verb!--doctype!   & \texttt{phd} / \texttt{msc} / \texttt{bsc} \\
  \verb!-s! / \verb!--docstatus! & \texttt{working} / \texttt{provisional} / \texttt{final} / \texttt{keep} \\
  \verb!-l! / \verb!--lang!      & Language code, e.g.\ \texttt{en}, \texttt{pt} \\
  \verb!--sdgs!                  & SDG numbers, e.g.\ \texttt{1,7,13} \\
  \verb!-c! / \verb!--cover!     & Cover-only build (fast preview) \\
  \verb!-f! / \verb!--force-school! & Apply school even in user mode \\
  \verb!-bdir! / \verb!--build-dir! & Workspace: omit=current, (empty)=temp, \texttt{-}=cached, path \\
  \verb!-k! / \verb!--keep-bdir! & Keep workspace after build \\
  \verb!-o! / \verb!--output-dir! & Destination for final PDF \\
  \verb!-r! / \verb!--rename-pdf! & Rename PDF to \texttt{univ-school-type-lang.pdf} \\
  \verb!-p! / \verb!--processor! & LaTeX engine: \texttt{lua} (default), \texttt{pdf}, \texttt{xe} \\
  \verb!--progress!              & \texttt{0}=silent, \texttt{1}=bar, \texttt{2}=realtime \\
  \verb!--lines!                 & Expected line count for progress bar \\
  \verb!--dry-run!               & Preview actions without building \\
  \verb!-v! / \verb!--verbose!   & Extra debug output \\
  \bottomrule
\end{ShadedTabular}
\end{center}
